\documentclass[11pt]{article}
\usepackage{mathblatt}


\begin{document}
\pfheftstil
\pfheftkopf{Grundwert $G$}{P10}
\pfrueckblick
\begin{pfaufg}{}{1.}{}
Ergänze die Tabelle Prozent | Bruch | Dezimalzahl.
\par\smallskip\noindent \begingroup\renewcommand{\arraystretch}{1.3}\begin{tabular}{|c|c|c|}\hline \textbf{Prozent} & \textbf{Bruch} & \textbf{Dezimalzahl} \\ \hline 1 \% & \rule{0pt}{3ex}\hspace{1.6cm} & \rule{0pt}{3ex}\hspace{1.6cm} \\ \hline 10 \% & \rule{0pt}{3ex}\hspace{1.6cm} & \rule{0pt}{3ex}\hspace{1.6cm} \\ \hline 20 \% & \rule{0pt}{3ex}\hspace{1.6cm} & \rule{0pt}{3ex}\hspace{1.6cm} \\ \hline 25 \% & \rule{0pt}{3ex}\hspace{1.6cm} & \rule{0pt}{3ex}\hspace{1.6cm} \\ \hline 50 \% & \rule{0pt}{3ex}\hspace{1.6cm} & \rule{0pt}{3ex}\hspace{1.6cm} \\ \hline 100 \% & \rule{0pt}{3ex}\hspace{1.6cm} & \rule{0pt}{3ex}\hspace{1.6cm} \\ \hline\end{tabular}\endgroup\par
\fusshilfe{\mbox{1:~0,01; 0,1; 0,2; 0,25; 0,5; 1}}
\end{pfaufg}
\begin{pfaufg}{}{2.}{}
Berechne 10 \% von 70 €.
\rechenplatz{1}
\fusshilfe{\mbox{2:~7 €}}
\end{pfaufg}
\pfstufekopf[sgrundwert]{Grundwert $G$}{in 2 der letzten 5 Prüfungen}
\pfgruppe{}
\begin{pfaufg}{}{3.}{}
$7\,\%$ eines Betrags sind $56\,€$. Rechne in der Tabelle: erst $1\,\%$, dann das Ganze ($100\,\%$).
\par\smallskip\noindent \begin{minipage}[t]{0.48\linewidth}\begingroup\renewcommand{\arraystretch}{1.35}\begin{tabular}[t]{|r|r|}\hline \textbf{Prozent} & \textbf{€} \\ \hline $7\,\%$ & $56$\,€ \\ \hline $1\,\%$ & \leerzelle \\ \hline $100\,\%$ & \leerzelle \\ \hline\end{tabular}\endgroup\hspace{2mm}{\small\color{mbgrau}\begin{tabular}[t]{@{}l@{}}\rule{0pt}{2.6ex}\\ $\downarrow : 7$\\ $\downarrow \cdot 100$\end{tabular}}\end{minipage}\hfill\begin{minipage}[t]{0.46\linewidth}\vspace{-2mm}\resizebox{\linewidth}{!}{\begin{minipage}{10.8cm}\streifen[0]{7}{$0$}{$100\,\%$}\end{minipage}}\end{minipage}\par
\fusshilfe{\mbox{3:~$800\,€$}}
\end{pfaufg}
\begin{pfaufg}{}{4.}{}
$16$ m sind $10\,\%$, $20\,\%$ oder $25\,\%$ eines Ganzen. Wie groß ist jeweils das Ganze?
\par\noindent \leerfeld[m] ; \leerfeld[m] ; \leerfeld[m] \par
\fusshilfe{\mbox{4:~$160$ m, $80$ m, $64$ m}}
\end{pfaufg}
\pfgruppe{}
\pfab{$G = W : p$}
\begin{pfaufg}{}{5.}{}
Ein Händler hat $24$ kg Äpfel verkauft. Das sind $10\,\%$ seiner Äpfel. Wie viel kg Äpfel hatte er?
\par\noindent \leerfeld[kg] \par
\fusshilfe{\mbox{5:~$240$ kg}}
\end{pfaufg}
\begin{pfaufg}{}{6.}{}
Im Tank sind noch $12$ l. Das sind $20\,\%$ der Tankfüllung. Wie viele Liter passen in den Tank?
\par\noindent \leerfeld[l] \par
\fusshilfe{\mbox{6:~$60$ l}}
\end{pfaufg}
\begin{pfaufg}{}{7.}{}
$21$ Kinder fahren mit dem Rad. Das sind $75\,\%$ der Klasse. Wie viele Kinder hat die Klasse?
\par\noindent \leerfeld \par
\fusshilfe{\mbox{7:~$28$ Kinder}}
\end{pfaufg}
\pfgruppe{}
\begin{pfaufg}{}{8.}{}
Was ist gesucht? Kreuze an. Rechne nicht.
\par\smallskip\noindent \begingroup\renewcommand{\arraystretch}{1.4}\begin{tabular}{|>{\raggedright\arraybackslash}p{4.3cm}|>{\centering\arraybackslash\hyphenpenalty=0}p{1.55cm}|>{\centering\arraybackslash\hyphenpenalty=0}p{1.55cm}|>{\centering\arraybackslash\hyphenpenalty=0}p{1.55cm}|>{\centering\arraybackslash\hyphenpenalty=0}p{1.55cm}|}\hline  & {\scriptsize\hspace{0pt}Prozentsatz $p$} & {\scriptsize\hspace{0pt}Prozentwert $W$} & {\scriptsize\hspace{0pt}Grundwert $G$} & {\scriptsize\hspace{0pt}neuer Wert} \\ \hline 2 von 16 Pfannkuchen sind mit Senf gefüllt. Wie viel Prozent sind das? & $\square$ & $\square$ & $\square$ & $\square$ \\ \hline Wie viel sind 30 \% von 70 €? & $\square$ & $\square$ & $\square$ & $\square$ \\ \hline 25 \% des Gewinns sind 200 €. Wie hoch ist der ganze Gewinn? & $\square$ & $\square$ & $\square$ & $\square$ \\ \hline 125 g + 20 \% mehr Inhalt & $\square$ & $\square$ & $\square$ & $\square$ \\ \hline Von 83 Elfmetern wurden 12 nicht verwandelt. Wie viel Prozent sind das? & $\square$ & $\square$ & $\square$ & $\square$ \\ \hline Ein Fahrrad kostet 550 €, bei Barzahlung gibt es 20 \% Rabatt. Wie viel Euro spart man? & $\square$ & $\square$ & $\square$ & $\square$ \\ \hline\end{tabular}\endgroup\par
\fusshilfe{\mbox{8:~1~$p$, 2~$W$, 3~$G$, 4~neuer Wert, 5~$p$, 6~$W$}}
\end{pfaufg}
\pfgruppe{}
\begin{pfaufg}{P10 ’23}{9.}{1 BE}
Mia gewinnt bei einem Wettbewerb Geld. Sie gibt 25 \% des Gewinns an ihre Eltern ab; das sind 200 €. Gib an, wie hoch der gesamte Gewinn war.
\rechenplatz{1}
\fusshilfe{\mbox{9:~800 €}, \mbox{Tipp: $G$ = 200 € : 0,25}}
\end{pfaufg}
\begin{pfbuendel}{gleiche Art \textperiodcentered{} 2$\times$ geprüft – kannst du überspringen}
\pfbpaar{\pfbglied{P10 ’25}{10.}{Ein T-Shirt ist um 6 € billiger geworden. Dieser Nachlass entspricht 20 \% des ursprünglichen Preises. Gib den ursprünglichen Preis an.}{1 BE}\fusshilfe{\mbox{10:~30 €}, \mbox{Tipp: $G$ = 6 € : 0,2}}}{}
\end{pfbuendel}
\begin{pfaufg}{NI ’26}{11.}{}
An einer Schule kommen 40 \% mit dem Fahrrad. Das sind 120 Kinder. Wie viele Kinder besuchen die Schule? Schreibe deinen Rechenweg auf.
\rechenplatz{1}
\fusshilfe{\mbox{11:~40 \% = 120 $\Rightarrow$ 10 \% = 30 $\Rightarrow$ 100 \% = 300 Kinder}}
\end{pfaufg}
\pfgruppe{}
\begin{pfaufg}{}{12.}{}
$9\,\%$ einer Menge sind $63$ kg. Wie viel ist $1\,\%$ der Menge? Rechne nicht weiter.
\par\noindent 1 \% = \leerfeld[kg] \par
\fusshilfe{\mbox{12:~$1\,\% = 7$ kg}}
\end{pfaufg}
\pfgruppe{}
\begin{pfaufg}{NI ’24}{13.}{}
Im Mai 2022 wurden in Norwegen 8 400 Elektro-Autos verkauft. Das waren 75 \% aller verkauften Autos. Berechne die Gesamtzahl.
\rechenplatz{1}
\fusshilfe{\mbox{13:~8 400 : 0,75 = 11 200}}
\end{pfaufg}
\pfgruppe{}
\begin{pfaufg}{BW ’23}{14.}{}
Am „Black Friday“ kostet ein Tablet 360,00 €. Der Preis wurde auf 80 \% gesenkt. Was kostete das Tablet vorher?
\fusshilfe{\mbox{14:~450 € (360 € sind 80 \%)}}
\end{pfaufg}
\begin{pfaufg}{}{15.}{}
$36\,\%$ einer Menge sind $81$ kg. Wie viel kg ist die ganze Menge?
\par\noindent \leerfeld[kg] \par
\fusshilfe{\mbox{15:~$225$ kg}}
\end{pfaufg}
\begin{pfaufg}{}{16.}{}
$64\,\%$ einer Strecke sind $52$ m. Wie lang ist die ganze Strecke?
\par\noindent \leerfeld[m] \par
\fusshilfe{\mbox{16:~$81{,}25$ m}}
\end{pfaufg}
\pfgruppe{}
\begin{pfaufg}{nach P10 ’19}{17.}{}
In einem Gefäß liegen 4 schwarze Kugeln. Sie sind $\tfrac{2}{3}$ aller Kugeln im Gefäß. Berechne, wie viele Kugeln insgesamt im Gefäß liegen.
\rechenplatz{1}
\fusshilfe{\mbox{17:~G = 4 : $\tfrac{2}{3}$ = 6 Kugeln ($\tfrac{1}{3}$ sind 2 Kugeln)}}
\end{pfaufg}
\pfgruppe{}
\begin{pfaufg}{nach P10 ’17}{18.}{}
In einem Becher steht der Quark 8 cm hoch; das sind 93 \% der Becherhöhe. Berechne die Höhe des Bechers.
\rechenplatz{1}
\fusshilfe{\mbox{18:~G = 8 cm : 0,93 $\approx$ 8,60 cm}}
\end{pfaufg}
\end{document}